\documentclass{article}
\usepackage[UTF8]{ctex}
\usepackage{amsmath,amssymb,amsfonts}
\usepackage{enumitem}
\usepackage{geometry}
\geometry{a4paper, margin=1in}
% === 错题本排版优化（不参与题目标识） ===
\setlength{\parindent}{0pt}
\setlength{\parskip}{0.35em plus 0.10em minus 0.08em}
\setlength{\emergencystretch}{3em}
\setlist{itemsep=0.25em, topsep=0.40em, parsep=0pt}
\allowdisplaybreaks[2]
% === 错题本排版优化结束 ===
\begin{document}
\section*{哈尔滨工业大学2026年高等代数真题}
\section*{76．哈尔滨工业大学2026年高等代数试卷}

\subsection{题目1}
一．求 7 次多项式 \(f(x)\) ，使得 \(f(x)+1\) 能被（x－1）整除，\(f(x)-1\) 能被（x＋1）整除。

\subsection{题目2}
二．计算行列式

\[
\left|\begin{array}{cccc}
2+a_{1}+b_{1} & a_{1}+b_{2} & \cdots & a_{1}+b_{n} \\
a_{2}+b_{1} & 2+a_{2}+b_{2} & \cdots & a_{2}+b_{n} \\
\vdots & \vdots & & \vdots \\
a_{n}+b_{1} & a_{n}+b_{2} & \cdots & 2+a_{n}+b_{n}
\end{array}\right| .
\]

\subsection{题目3}
三．设 \(A\) 是实数域上的 \(m \times n\) 矩阵，\(X\) 是 \(n \times m\) 矩阵，\(X\) 的元素由独立的未知数构成，\(m \leq n\) 。证明：

\[
A X=E_{m}
\]

有解的充分必要条件是秩 \((A)=m\) ．

\subsection{题目4}
四．求平方为零的全体三阶实方阵．

\subsection{题目5}
五．设数域 \(P\) 上的 \(n\) 元二次型 \(f\left(x_{1}, x_{2}, \cdots, x_{n}\right)\) 对应的方阵 \(S\) 的顺序主子式 \(d_{1}, d_{2}, \cdots, d_{n}\) 均不为 0 ，其中 \(d_{i}\) 是 \(i\) 阶顺序主子式。求 \(f\left(x_{1}, x_{2}, \cdots, x_{n}\right)\) 的标准形。

\subsection{题目6}
六．设矩阵 \(A, B\) 满足 \(A B=2 A+3 B\) ，其中

\[
A=\left(\begin{array}{ccc}
1 & 1 & 1 \\
1 & 0 & -1 \\
1 & -1 & -3
\end{array}\right)
\]

求 \(B\) 的特征值．

\begin{itemize}
\end{itemize}

\subsection{题目7}
原题边界“七．”：判断如下命题是否成立，并说明理由：存在不全为零的有理数 \(a, b, c, d\) ，使得
\[
x^{3}+6 x+1 \mid a+b x^{2}+c x^{4}+d x^{6}
\]

\subsection{题目8}
原题边界“八．”：设实数 \(a, b, c\) 满足 \(a+2 b+2 c=1\) ，求 \(a^{2}+b^{2}+c^{2}\) 的取值范围．

\subsection{题目7}
九．设实数域上方阵 \(A=\left(a_{i j}\right)_{8 \times 8}\) 满足 \(\left|a_{i j}\right| \leq 1\) ，其中 \(i, j=1,2, \cdots, 8\) ，求 \(|A|\) 的最大值．

\subsection{题目8}
十．判断如下命题是否成立？并说明理由：若实数域上方阵 \(A\) 的行向量两两正交，则其列向量也两两正交。
\end{document}
